% ═══════════════════════════════════════════════════════════════
% MT-03d: Minitest Wiederholung Mi casa
% Spanisch Klasse 7 | Sequenz 3: Mi casa
% ═══════════════════════════════════════════════════════════════
% Dauer: 20 Minuten
% Punkte: 32 BE
% Inhalt: Räume, hay + Möbel, Farbadjektive, Zimmer beschreiben
% ═══════════════════════════════════════════════════════════════

\documentclass[11pt, a4paper]{article}

\usepackage[utf8]{inputenc}
\usepackage[T1]{fontenc}
\usepackage[ngerman]{babel}
\usepackage{helvet}
\renewcommand{\familydefault}{\sfdefault}

\usepackage[margin=2cm]{geometry}
\usepackage{enumitem}
\usepackage{tabularx}
\usepackage{booktabs}

\usepackage{xcolor}
\usepackage{tcolorbox}
\tcbuselibrary{skins}

\usepackage{fancyhdr}
\usepackage{lastpage}
\usepackage{needspace}

% FARBEN
\definecolor{spanisch-color}{HTML}{c41e3a}
\definecolor{grau-color}{HTML}{888888}
\definecolor{hellgrau-color}{HTML}{f5f5f5}

\colorlet{spanisch}{spanisch-color}
\colorlet{grau}{grau-color}
\colorlet{hellgrau}{hellgrau-color}

\newcommand{\fachfarbe}{spanisch}

\newcommand{\thema}{Minitest: Wiederholung Mi casa}
\newcommand{\fach}{Spanisch}
\newcommand{\klasse}{7}
\newcommand{\maxpunkte}{32}
\newcommand{\zeit}{20 Minuten}
\newcommand{\datum}{\today}

\pagestyle{fancy}
\fancyhf{}
\renewcommand{\headrulewidth}{0.4pt}
\renewcommand{\footrulewidth}{0.4pt}

\lhead{\fach{} -- Klasse \klasse}
\chead{\textbf{MT-03d}}
\rhead{\datum}

\lfoot{\zeit}
\cfoot{}
\rfoot{Seite \thepage\ von \pageref{LastPage}}

\newcommand{\punktebox}[1]{\hfill\fbox{\small \_/#1}}
\newcommand{\luecke}[1]{\rule{#1}{0.4pt}}

\newtcolorbox{infobox}{
    colback=hellgrau,
    colframe=\fachfarbe,
    boxrule=1pt,
    arc=0pt,
    left=8pt, right=8pt, top=8pt, bottom=8pt
}

\newtcolorbox{hintbox}{
    colback=white,
    colframe=blue!50!black,
    boxrule=1pt,
    arc=0pt,
    left=8pt, right=8pt, top=8pt, bottom=8pt
}

\newenvironment{aufgabe}[1]{%
    \needspace{5\baselineskip}%
    \nopagebreak[4]%
    \vspace{10pt}
    \noindent\textbf{\textcolor{\fachfarbe}{Aufgabe #1}}
    \nopagebreak[4]%
    \vspace{5pt}
    \nopagebreak[4]%
    \begin{list}{}{\leftmargin=0pt}
    \item[]
}{%
    \end{list}
}

\newcommand{\es}[1]{\textcolor{\fachfarbe}{\textbf{#1}}}

\begin{document}

% ═══════════════════════════════════════════════════════════════
% KOPFBEREICH
% ═══════════════════════════════════════════════════════════════

\begin{center}
    {\Large\bfseries\textcolor{\fachfarbe}{\thema}}
\end{center}

\vspace{5pt}

\noindent
\begin{tabularx}{\textwidth}{|X|X|X|}
\hline
\textbf{Name:} \luecke{4cm} &
\textbf{Klasse:} \klasse &
\textbf{Datum:} \luecke{2cm} \\
\hline
\end{tabularx}

\vspace{10pt}

\begin{infobox}
\textbf{Zeit:} \zeit{} | \textbf{Punkte:} \maxpunkte{} BE | \textbf{Hilfsmittel:} keine
\end{infobox}

\vspace{15pt}

% ═══════════════════════════════════════════════════════════════
% AUFGABE 1: Räume und Artikel (6 BE)
% ═══════════════════════════════════════════════════════════════

\begin{aufgabe}{1 -- Räume und Artikel}
Setze den richtigen Artikel ein: \es{el} oder \es{la}? \punktebox{6}
\end{aufgabe}

\noindent
a) \luecke{1.5cm} cocina \hfill
b) \luecke{1.5cm} salón \\[0.8em]
c) \luecke{1.5cm} dormitorio \hfill
d) \luecke{1.5cm} baño \\[0.8em]
e) \luecke{1.5cm} terraza \hfill
f) \luecke{1.5cm} jardín \\

\vspace{15pt}

% ═══════════════════════════════════════════════════════════════
% AUFGABE 2: hay + Möbel (8 BE)
% ═══════════════════════════════════════════════════════════════

\begin{aufgabe}{2 -- hay + Möbel}
Ergänze mit \es{hay} und dem passenden Möbelstück (Artikel nicht vergessen!). \punktebox{8}
\end{aufgabe}

\noindent
a) En el dormitorio \luecke{5cm}. \hfill \textit{(Bett)}

\vspace{0.8em}

\noindent
b) En el salón \luecke{5cm}. \hfill \textit{(Sofa)}

\vspace{0.8em}

\noindent
c) En la cocina \luecke{5cm}. \hfill \textit{(Tisch)}

\vspace{0.8em}

\noindent
d) En mi habitación \luecke{5cm}. \hfill \textit{(Schrank)}

\vspace{15pt}

% ═══════════════════════════════════════════════════════════════
% AUFGABE 3: Farbadjektive (8 BE)
% ═══════════════════════════════════════════════════════════════

\begin{aufgabe}{3 -- Farbadjektive anpassen}
Setze das Farbadjektiv in der richtigen Form ein. \punktebox{8}
\end{aufgabe}

\noindent
a) La mesa es \luecke{3cm}. \hfill \textit{(rot)}

\vspace{0.8em}

\noindent
b) El sofá es \luecke{3cm}. \hfill \textit{(grün)}

\vspace{0.8em}

\noindent
c) Las sillas son \luecke{3cm}. \hfill \textit{(weiß)}

\vspace{0.8em}

\noindent
d) Los armarios son \luecke{3cm}. \hfill \textit{(schwarz)}

\vspace{15pt}

% ═══════════════════════════════════════════════════════════════
% AUFGABE 4: Zimmer beschreiben (10 BE)
% ═══════════════════════════════════════════════════════════════

\begin{aufgabe}{4 -- Dein Zimmer beschreiben}
Beschreibe dein Zimmer auf Spanisch. Schreibe mindestens \textbf{5 Sätze}. \punktebox{10}
\end{aufgabe}

\begin{hintbox}
\textbf{Hilfe:} Verwende \es{hay}, \es{es}, Farbadjektive, Möbel und Räume. \\
Beispiel: \textit{En mi habitación hay una cama. La cama es azul.}
\end{hintbox}

\vspace{0.8em}

\noindent
\rule{\textwidth}{0.4pt}

\vspace{1.2em}

\noindent
\rule{\textwidth}{0.4pt}

\vspace{1.2em}

\noindent
\rule{\textwidth}{0.4pt}

\vspace{1.2em}

\noindent
\rule{\textwidth}{0.4pt}

\vspace{1.2em}

\noindent
\rule{\textwidth}{0.4pt}

\vspace{1.2em}

\noindent
\rule{\textwidth}{0.4pt}

\vspace{0.5em}

\noindent\textit{(2 BE pro korrektem Satz, max. 5 Sätze = 10 BE)}

% ═══════════════════════════════════════════════════════════════
% PUNKTEÜBERSICHT
% ═══════════════════════════════════════════════════════════════

\vspace{25pt}
\hrule
\vspace{10pt}

\noindent
\begin{tabularx}{\textwidth}{|X|c|c|c|c||c|}
\hline
\textbf{Aufgabe} & 1 & 2 & 3 & 4 & \textbf{Gesamt} \\
\hline
\textbf{max. Punkte} & 6 & 8 & 8 & 10 & \textbf{32} \\
\hline
\textbf{erreicht} & & & & & \\[0.8em]
\hline
\end{tabularx}

\vspace{10pt}

\begin{flushright}
\textbf{Note:} \luecke{2cm}
\end{flushright}

\end{document}
